\documentclass[../../../include/open-logic-section]{subfiles}

\begin{document}

\olfileid[id]{his}{bio}{gen}

\olsection{Gerhard Gentzen}

\olphoto{gentzen-gerhard}{Gerhard Gentzen}

Gerhard Gentzen terutama dikenal sebagai pencipta teori bukti struktural,
khususnya sistem !!{derivation} kalkulus sekuen dan deduksi alami. Ia lahir
pada 24 November 1909 di Greifswald, Jerman. Gerhard bersekolah di rumah
selama tiga tahun sebelum memasuki sekolah persiapan, tempat tingkat
pendidikannya tertinggal dari kebanyakan teman sekelasnya. Meskipun demikian,
ia merupakan murid yang cemerlang dan menunjukkan bakat kuat dalam
matematika. Minatnya beragam; misalnya, ia juga menulis puisi untuk ibunya
dan naskah drama bagi teater sekolah.

Gentzen memulai studi universitasnya di Universitas Greifswald, tetapi
kemudian berpindah-pindah ke G\"{o}ttingen, München, dan Berlin. Ia meraih
gelar doktor pada tahun 1933 dari Universitas G\"{o}ttingen di bawah
bimbingan Hermann Weyl. (Paul Bernays membimbing sebagian besar pekerjaannya,
tetapi diberhentikan dari universitas oleh Nazi.) Pada tahun 1934, Gentzen
mulai bekerja sebagai asisten David Hilbert. Pada tahun yang sama ia
mengembangkan sistem !!{derivation} kalkulus sekuen dan deduksi alami dalam
makalahnya \emph{Untersuchungen \"{u}ber das logische Schlie\ss en I--II
  [Penyelidikan tentang Deduksi Logis I--II]}. Ia membuktikan konsistensi
aksioma Peano pada tahun 1936.

Hubungan Gentzen dengan Nazi bersifat rumit. Pada saat pembimbingnya, Bernays,
dipaksa meninggalkan Jerman, Gentzen bergabung dengan cabang SA di
universitas, organisasi paramiliter Nazi. Seperti banyak orang Jerman, ia
menjadi anggota Partai Nazi. Selama perang, ia bertugas sebagai perwira
telekomunikasi bagi unit intelijen udara. Namun, pada tahun 1942 ia
dibebastugaskan karena mengalami gangguan saraf. Tidak jelas apakah kesetiaan
Gentzen berada pada Partai Nazi atau apakah ia bergabung demi menjamin
keberhasilan akademiknya.

Pada tahun 1943, Gentzen ditawari jabatan akademik di Institut Matematika
Universitas Jerman di Praha, dan ia menerimanya. Namun, pada tahun 1945 warga
Praha memberontak melawan pendudukan Jerman. Pasukan Soviet tiba di kota itu
dan menangkap semua profesor di universitas tersebut. Karena keanggotaannya
dalam organisasi Nazi, Gentzen dibawa ke kamp kerja paksa. Ia meninggal
akibat malnutrisi di selnya pada 4 Agustus 1945 dalam usia 35 tahun.

\begin{reading}
Untuk biografi lengkap Gentzen, lihat \citet{Menzler-Trott2007}. Bacaan
menarik mengenai para matematikawan di bawah kekuasaan Nazi, yang juga
memberikan catatan singkat tentang kehidupan Gentzen, disajikan oleh
\citet{Segal2014}. Makalah Gentzen tentang deduksi logis tersedia dalam
bahasa Jerman aslinya \citep{Gentzen1935a,Gentzen1935b}. Terjemahan bahasa
Inggris makalah-makalah Gentzen telah dihimpun dalam satu jilid oleh
\citet{Gentzen1969}, yang juga memuat sketsa biografis.
\end{reading}

\end{document}
